\color{unidarkgreen}{\subsection{Комбинированный пусковой орган напряжения (КПОН)}}
\color{black}

\begin{enumerate}[label=\arabic{section}.\arabic{subsection}.\arabic*,  align=left, labelsep=4pt, leftmargin=0pt, itemindent=57pt]

\item
Комбинированный пусковой орган по напряжению низкой и средней сторон трансформатора используется для повышения чувствительности МТЗ ВН и МТЗ СН. 

\item
Алгоритм работы КПОН приведен на рисунке \ref{pict:KPON}. Уставки КПОН приведены в таблице \ref{app_set:kpon}.

\vspace{2mm}
\begin{figure}[H]
\centering
\includegraphics[width=0.9\textwidth,height=0.9\textheight,keepaspectratio]{\fbpath/07.01 РЗТ/081.0371 КПОН/img/KPON.pdf}
\captionof{figure}{Функциональная схема логики КПОН}\label{pict:KPON}
\end{figure}

\item
Цепи напряжения подключаются к ТН, установленным на шинах низшего и среднего напряжения.

\item
Внутренний комбинированный пусковой орган напряжения (КПОН), контролирует минимальное линейное напряжение и максимальное напряжение обратной последовательности низкой и средней стороны трансформатора.  

\item
Порог срабатывания пуска по минимальному линейному напряжению задается уставками \textbf{«Uл НН1(НН) 1СШ»}, \textbf{«Uл НН1(НН) 2СШ»}, \textbf{«Uл НН2(СН) 1СШ»} и \textbf{«Uл НН2(СН) 2СШ»}, порог срабатывания пуска по максимальному напряжению обратной последовательности задается уставками \textbf{«U2 НН1(НН) 1СШ»}, \textbf{«U2 НН1(НН) 2СШ»}, \textbf{«U2 НН2(СН) 1СШ»} и \textbf{«U2 НН2(СН) 2СШ»}. 

\item
Ввод в работу КПОН осуществляется независимо для каждой стороны с помощью программных переключателей \textbf{«КПОН НН1(НН) 1СШ»}, \textbf{«КПОН НН1(НН) 2СШ»}, \textbf{«КПОН НН2(СН) 1СШ»} и \textbf{«КПОН НН2 СН) 2СШ»}.

\item
Оперативный ввод/вывод КПОН производится с помощью БУ (блоков управления функциями) \textit{«Вывод КПОН НН1(НН) 1СШ»}, \textit{«Вывод КПОН НН1(НН) 2СШ»}, \textit{«Вывод КПОН НН2(СН) 1СШ»} и \textit{«Вывод КПОН НН2(СН) 2СШ»}.

\item
Выбор типа пуска по напряжению соответствующей стороны задается программными переключателями \textbf{«Тип КПОН НН1(НН) 1СШ»}, \textbf{«Тип КПОН НН1(НН) 2СШ»}, \textbf{«Тип КПОН НН2(СН) 1СШ»} и \textbf{«Тип КПОН НН2(СН) 2СШ»}, имеющими три положения:
\begin{itemize}
\item \textit{«Umin, U2»} -- пуск по минимальному линейному напряжению или максимальному напряжению обратной последовательности;
\item \textit{«Umin»} -- пуск по минимальному линейному напряжению;
\item \textit{«Внеш»} -- пуск по внешнему входному сигналу «Внеш. КПОН» соответствующей стороны. 
\end{itemize}

\item
Предусмотрена возможность контроля неисправности цепей напряжения средней и низкой сторон трансформатора по внутренней логике КЦН. При появлении неисправности цепей напряжения, действие логики КПОН соответствующей стороны задается программным переключателями «Неиспр ТН НН2(СН) 1СШ», «Неиспр ТН НН2(СН) 2СШ», «Неиспр ТН НН1(НН) 1СШ» и «Неиспр ТН НН1(НН) 2СШ», имеющими три положения:
\begin{itemize}
\item \textit{«Выведено»} -- режим работы КПОН не зависит от КЦН;
\item \textit{«Пуск КПОН»} -- при появлении неисправности в цепях ТН, КПОН всегда разрешен;
\item \textit{«Блок КПОН»} -- при появлении неисправности в цепях ТН, КПОН выводится из работы. 
\end{itemize}












































\end{enumerate} 